\documentclass[a4paper]{article}
\usepackage[pdftex]{graphicx}
\usepackage[utf8]{inputenc}
\usepackage{enumerate}
\usepackage{href-ul}
\hypersetup{
	colorlinks=true,
	linkcolor=blue,
	urlcolor=blue}
\usepackage{siunitx}
\sisetup{locale=DE} 
\usepackage{icomma}
\usepackage{amssymb}
\usepackage{geometry}
\geometry{a4paper, top=15mm, left=15mm, right=15mm, bottom=15mm,
	headsep=10mm, footskip=12mm}

\begin{document}

\begin{enumerate}[1.]
\item {\bf Masse und Trägheit}

\begin{enumerate}[a)]
\item Unter welcher Voraussetzung verhält sich ein Körper träge? Was versteht man unter einem trägen Verhalten?
\item Obwohl Thomas auf ebener Strecke mit all seiner Kraft in die Pedale tritt, fährt er mit konstanter Geschwindigkeit. Erläutere, woran das liegt.
\item Schnee an den Schuhen kann man durch kräftiges Stampfen auf den Boden entfernen. Erkläre diesen Sachverhalt aus physikalischer Sicht.
\end{enumerate}


\item {\bf Hebel; Maschinen} \\
Ein symmetrisch gelagerter Hebel wird mit einer Seilwinde über eine Seil-Rollen-Kombination bewegt. Das angehängte Gewicht G beträgt $F_G=\SI{1,50}{\kilo\newton}$. Am Hebel gilt: $a_1=\SI{1,50}{\meter}$ und $a_2=\SI{2,50}{\meter}$.
\begin{enumerate}[a)]
\item Mit welcher Kraft $F_S$ muss die Seilwinde am Seil ziehen, dass Gleichgewicht herrscht ? 
\item Die Seilwinde besteht aus einer Kurbel mit dem Radius $ R=\SI{40,0}{\centi\meter}$ und der Seiltrommel mit dem Radius $ r=\SI{10,0}{\centi\meter}$ . \\
Berechne, welche Kraft $F_K$ in diesem Fall an der Kurbel wirken muss
\item Die Seilverankerung im Punkt A hält höchstens 600 N aus. Wie weit kann sich der Angriffspunkt des Gewichtes G am linken Hebelarm unter dieser Bedingung höchstens von der Hebelachse entfernen?
\end{enumerate}

\includegraphics[width=12 cm]{hebel119.pdf}

\item {\bf Drehmoment}

\begin{enumerate}[a)]
\item Der Griff einer Beißzange wird $a_1=\SI{15}{\centi\meter}$ vom Drehpunkt entfernt mit einer Kraft von 
$F_1=\SI{180}{\newton}$ zusammengedrückt, um einen Nagel, der $a_2=\SI{3,5}{\centi\meter}$ vom Drehpunkt entfernt ist, durchzutrennen.\\
Welche Kraft wirkt auf den Nagel ?
\item Senkrecht zu einer symmetrisch gelagerten Hebelstange wirken auf der linken Seite $a_1=\SI{3,0}{\centi\meter}$ vom Drehpunkt entfernt eine Kraft von $F_1=\SI{6,0}{\newton}$ und auf der rechten Seite $a_2=\SI{8,0}{\centi\meter}$ vom Drehpunkt entfernt eine Kraft von $F_2=\SI{3,0}{\newton}$ \\
Fertige eine Skizze an und ergänze eine dritte Kraft so, dass Gleichgewicht herrscht. (Begründung!) 
\end{enumerate}

\item {\bf Flaschenzug}

\begin{minipage}{0.45\textwidth}
\begin{enumerate}[a)]
\item Wie groß muss die Kraft im Punkt C sein, dass Gleichgewicht herrscht ?
\item Wie groß sind die Kräfte in den Aufhängepunkten  A und B?
\end{enumerate}
\vfill
\end{minipage}
\hfill
\begin{minipage}{0.5\textwidth}
\includegraphics[width=5 cm]{flaschzug121.pdf}
\end{minipage}

\end{enumerate} 

\fbox{
	\begin{minipage}{0.5\textwidth}
		Zur Lösung bitte \href{https://www.okuyakl.de/phys/p8mecL406/ll406.pdf}{hier klicken} oder den QR-Code scannen.\\
	Weitere Arbeitsblätter gibt es unter 
	
	\href{https://www.okuyakl.de}{www.okuyakl.de}
	\end{minipage}
	\hfill
	\begin{minipage}{0.4\textwidth}
		\includegraphics[width=1.5 cm]{../../viecher/zwe03}
		\includegraphics[width=3 cm]{qr406}
		\includegraphics[width=2 cm]{../../viecher/afanticon1}
		
	\end{minipage}}

\end{document}%Lösung-------------------------------------------
